\documentclass[11pt]{article}

\usepackage[margin=1in]{geometry}
\usepackage{amsmath,amssymb,amsthm,mathtools,microtype}
\usepackage[T1]{fontenc}
\usepackage[hidelinks]{hyperref}

\newtheorem{theorem}{Theorem}[section]
\newtheorem{remark}[theorem]{Remark}

\newcommand{\C}{\mathbb C}
\newcommand{\N}{\mathbb N}
\newcommand{\E}{\mathbb E}
\newcommand{\Supp}{\operatorname{Supp}}

\title{Sparse Minimality of Gaussian-Moments Counterexamples\\
in Dimension Three}
\author{Roy van Rijn\\{\small Independent researcher}}
\date{24 July 2026}
\hypersetup{
  pdftitle={Sparse Minimality of Gaussian-Moments Counterexamples in Dimension Three},
  pdfauthor={Roy van Rijn},
  pdfsubject={Sparse minimal counterexamples to the Gaussian Moments Conjecture in dimension three},
  pdfkeywords={Gaussian Moments Conjecture, Gaussian moments,
    sparse polynomials, Groebner bases, computer-assisted proof}
}

\begin{document}
\maketitle

\begin{abstract}
The Gaussian Moments Conjecture fails in every dimension at least three.
We study minimal failure in the first failing dimension.  An exhaustive
exact search over all \(59{,}535\) supports having degree at most four and
at most four monomials shows that every pure-moment-null polynomial in this
range still satisfies the conjectural mixed-moment conclusion.  Thus
Long's five-term quartic is Pareto-minimal with respect to degree and
circular-coordinate support size.  The theorem is computer-assisted:
moment ideals through order twenty-four are computed directly in Laurent
coordinate rings over \(\mathbb Q\), and the three genuine surviving
families are proved harmless at all orders.
\end{abstract}

\section{Introduction}

Let \(X_1,\ldots,X_d\) be independent standard real Gaussian random
variables.  The Gaussian Moments Conjecture, introduced by Derksen, van den
Essen, and Zhao~\cite{DVZ}, asserts that if
\[
  P\in\C[X_1,\ldots,X_d],
  \qquad \E(P^m)=0\quad(m\geq1),
\]
then, for every \(Q\in\C[X_1,\ldots,X_d]\),
\[
  \E(QP^m)=0
  \qquad\text{for all sufficiently large }m.
\]
Long recently found explicit counterexamples in every dimension
\(d\geq3\)~\cite{Long}.  His three-variable witness, in circular
coordinates \(W,Z\) and an independent real Gaussian \(T\), is
\begin{equation}\label{eq:long}
 P_{\mathrm L}
 =W+WZ-T^2-\frac32ZT^2-\frac12Z^2T^2.
\end{equation}
It has degree four and five monomials, all its pure moments vanish, and
\(\E(ZP_{\mathrm L}^m)=m!\) for every \(m\geq1\).  We place this witness at
the first corner of the sparse degree--support plane.

\begin{theorem}[Sparse Pareto minimality]\label{thm:sparse}
Let \(W,Z,T\) be Gaussian Witt coordinates, so that
\[
 \E(W^iZ^jT^k)=
 \begin{cases}
   i!(k-1)!!,&i=j\text{ and }k\text{ is even},\\
   0,&\text{otherwise}.
 \end{cases}
\]
Suppose that \(P\in\C[W,Z,T]\) has total degree at most four and at most four
nonzero monomials.  If \(\E(P^m)=0\) for every \(m\geq1\), then, for every
\(Q\in\C[W,Z,T]\),
\[
 \E(QP^m)=0
 \qquad\text{for all sufficiently large }m.
\]
Consequently no GMC counterexample has both degree at most four and support
size at most four in a Gaussian Witt frame.
\end{theorem}

In particular, \eqref{eq:long} is minimal for the product partial order on
\((\deg P,\#\Supp(P))\).  The theorem does not separately prove that degree
four is globally minimal: a denser cubic is not excluded.  Nor does it
exclude a four-term counterexample of degree greater than four.

\section{Moment ideals}

The finite part of the proof is an exhaustive exact coefficient-scheme
calculation; the all-order interpretation of its survivors is conceptual.
Let
\[
 \mathcal M_4
 =\{W^iZ^jT^k:i,j,k\geq0,\ i+j+k\leq4\}.
\]
This set has \(35\) elements.  Fix a support \(S\subset\mathcal M_4\), write
\[
 P_S=\sum_{s\in S}c_ss,
\]
and work on the coefficient torus \(\prod_{s\in S}c_s\neq0\).  For
\(s=(i_s,j_s,k_s)\) and a multiplicity vector
\(\nu=(\nu_s)_{s\in S}\in\mathbb N^S\), put
\[
 |\nu|=\sum_s\nu_s,\quad
 I(\nu)=\sum_s\nu_si_s,\quad
 J(\nu)=\sum_s\nu_sj_s,\quad
K(\nu)=\sum_s\nu_sk_s.
\]
Before expanding any Wick coefficients, associate to \(S\) its contraction
semigroup
\[
 \Gamma_S=
 \left\{\nu\in\mathbb N^S:
   \sum_s\nu_s(i_s-j_s)=0,\quad
   \sum_s\nu_sk_s\equiv0\pmod2
 \right\}.
\]
Its Hilbert basis \(H_S\) consists of the primitive contractions.  This
basis is computed exactly in two stages.  First ignore parity.  Pottier's
bound~\cite{Pottier} shows that every Hilbert-basis vector for the one-row
charge equation has
\[
  |\nu|\leq1+\sum_{s\in S}|i_s-j_s|.
\]
The even-parity subsemigroup is then generated by the even charge atoms and
the pairwise sums, with repetition, of odd charge atoms; removing
decomposable candidates gives \(H_S\).

Wick contraction gives the exact moment polynomial
\begin{equation}\label{eq:sparse-moment}
 \E(P_S^m)
 =
 \sum_{\substack{|\nu|=m,\ I(\nu)=J(\nu)\\K(\nu)\ {\rm even}}}
 \frac{m!}{\prod_s\nu_s!}\,
 I(\nu)!\,(K(\nu)-1)!!\prod_sc_s^{\nu_s}.
\end{equation}
Thus the truncated full-support coefficient scheme is most naturally
defined directly on the coefficient torus by the Laurent ideal
\begin{equation}\label{eq:laurent-ideal}
 J_{S,M}
 =
 \left\langle\E(P_S),\ldots,\E(P_S^M)\right\rangle
 \subseteq \mathbb Q[c_s^{\pm1}:s\in S].
\end{equation}
The equations are homogeneous, so one coefficient may be normalized to
one.  There is a second coefficient symmetry.  The
Gaussian-preserving torus
\[
  W\longmapsto tW,\qquad Z\longmapsto t^{-1}Z,\qquad T\longmapsto T
\]
acts on the coefficient of \(s=W^{i_s}Z^{j_s}T^{k_s}\) with charge
\(q_s=i_s-j_s\).  Thus the coefficient torus carries the action
\[
  (\lambda,t)\cdot c_s=\lambda t^{q_s}c_s,
\]
and every moment equation is invariant under \(t\) and homogeneous under
\(\lambda\).

If a support contains coefficients \(c_a,c_b\) of distinct charges, then
over an algebraic closure every orbit meets the slice \(c_a=c_b=1\):
one solves
\[
  t^{q_b-q_a}=c_a/c_b,\qquad
  \lambda=(t^{q_a}c_a)^{-1}.
\]
Equivalently, with \(d=q_b-q_a\), quotient coordinates are the invariant
Laurent ratios
\[
  y_s=c_s^d c_a^{q_s-q_b}c_b^{q_a-q_s}
  \qquad(s\neq a,b).
\]
On the slice they become \(y_s=c_s^d\), so the slice is a finite cover of
the invariant-ratio quotient.  Consequently the four-term elimination has
only two active coefficient variables whenever two charges occur.  If all
charges agree, the torus acts through overall scaling; such an active
non-one-sided support necessarily has charge zero, and the calculation
falls back to one coefficient normalization.

There are
\[
 \sum_{r=1}^4\binom{35}{r}=59{,}535
\]
raw supports.  We identify a support with its reflection under
\(W\leftrightarrow Z\).  Two elementary filters precede the Groebner
calculation.  If \(H_S\) has one element, then \(\Gamma_S\) has rank one
and the moment at the degree of its generator consists of a single
nonzero coefficient monomial; the coefficient torus contains no solution.
More generally, the contraction vectors are enumerated by total degree
from \(H_S\), and a degree containing exactly one vector gives the same
rejection before any factorial Wick weights are expanded.  Supports with
the same ordered charge/parity signature share this enumeration.  Only
the remaining supports are expanded using \eqref{eq:sparse-moment}.  If no
contraction occurs, all charges \(i_s-j_s\) have one strict sign and the
mixed-moment conclusion follows by bounded weight.  Possible additional
factorizations among disjoint primitive contractions are not used in the
certificate.

For completeness, the latter test is finite in this range.  A zero charge
contracts in one or two copies, according to its \(T\)-parity.  Opposite
nonzero charges \(q>0>r\), with \(|q|,|r|\leq4\), have a primitive
zero-sum relation of length
\[
 \frac{q+|r|}{\gcd(q,|r|)}\leq8;
\]
doubling corrects odd \(T\)-parity.  Hence every support which is not
strictly one-sided has a contraction by order sixteen.

\section{Exact enumeration}

After the preceding filters, the four-term stratum contains \(5{,}718\)
coefficient-torus ideals requiring elimination.  Each is computed in the
Laurent ring of the two-normalization slice when its charge data have rank
two, with the one-normalization fallback just described.  No inverse or
Rabinowitsch variable is adjoined: extension to the Laurent ring and
contraction are performed by ideal quotients in the genuine quotient
coordinates.
\begin{samepage}
Exact Groebner bases over \(\mathbb Q\) for \(J_{S,20}\) leave precisely
the supports
\[
\begin{split}
 &\{Z,ZT^2,WZ^2,W^3T\},\\
 &\{ZT,Z^2,WT,W^2\},\\
 &\{ZT,ZT^3,WZ^2T,W^3\}.
\end{split}
\]
\end{samepage}
The first and third acquire the unit ideal after adjoining moments through
order twenty-four.  The same exact search in the one-, two-, and three-term
strata leaves no active one- or two-term family and two three-term
families.  Including the four-term survivor, the genuine active null
families are
\begin{align}
 P_{\rm lin}&=T+aZ+bW,&1+2ab&=0,                 \label{eq:null-linear}\\
 P_{\rm sh}&=Z+aWT-\frac12a^2W^3,                \label{eq:null-shear}\\
 P_{\rm pr}&=ZT+aZ^2+\frac1{2a^2}WT
                    -\frac1{4a^3}W^2,\quad a\neq0. \label{eq:null-product}
\end{align}
For the three survivors only, clearing Laurent monomials exports triangular
lexicographic certificates.  With the first coefficient normalized to one,
their nontrivial equations are respectively
\[
 1+2c_1c_2,\qquad
 2c_2+c_1^2,\qquad
 2c_1^2c_2-1,\quad c_3+c_1c_2^2.
\]
These formulas both classify the coefficient schemes and make the
survivors easy to interpret.

\section{All-order interpretation of the survivors}

The family \eqref{eq:null-linear} is an isotropic linear form.  A complex
orthogonal change takes it to a scalar multiple of \(Z\).  A fixed
polynomial \(Q\) has bounded \(W\)-degree, so
\(\E(QZ^m)=0\) for all sufficiently large \(m\).

For \eqref{eq:null-shear}, define
\[
 D=W\partial_T-T\partial_Z.
\]
Then \(D(2WZ+T^2)=0\), \(D(W)=0\), and
\[
 P_{\rm sh}=\exp(-aWD)(Z).
\]
Gaussian integration by parts gives \(\E(DF)=0\) for every polynomial
\(F\), and therefore also \(\E(WDF)=0\).  Since \(WD\) is locally nilpotent,
\(\phi=\exp(-aWD)\) is a polynomial automorphism preserving the Gaussian
functional.  For fixed \(Q\),
\[
 \E(QP_{\rm sh}^m)
 =\E\bigl(\phi^{-1}(Q)Z^m\bigr),
\]
which vanishes eventually by the same bounded-weight argument.

Finally,
\begin{equation}\label{eq:product-factorization}
 P_{\rm pr}
 =
 \left(Z+\frac1{2a^2}W\right)
 \left(T+aZ-\frac1{2a}W\right).
\end{equation}
The second factor is isotropic and orthogonal to the first.  After an
orthogonal Witt change, \eqref{eq:product-factorization} is a scalar
multiple of \(TZ\).  Every \(P_{\rm pr}^m\) then has isotropic weight \(m\),
which a fixed \(Q\) can balance only finitely often.  All three families
satisfy GMC.  Together with the strictly one-sided supports, this proves
Theorem~\ref{thm:sparse}.

\section{Reproducibility and scope}

The enumeration is implemented with exact integer Wick coefficients and
Singular's rational Groebner bases.  It computes each contraction Hilbert
basis before moment expansion, caches contraction vectors by charge/parity
signature, and enumerates supports rather than reading a stored survivor
list.  It performs every emptiness test over \(\mathbb Q\) in the Laurent
ring of the coefficient-symmetry quotient, without auxiliary saturation
variables.  Polynomial denominators are cleared only to export the three
displayed triangular certificates.  The replay also verifies the
invariant-ratio exponents and checks \eqref{eq:product-factorization}.  It
is run from the repository root by
\begin{verbatim}
.venv/bin/python scripts/verify_three_real_gmc_sparse_pareto.py
\end{verbatim}
The complete replay takes under a minute on a laptop.  Modular runs are
used only as optional discovery filters and are not part of the exact
certificate.

\section*{Acknowledgments and declaration of generative AI use}

The author used OpenAI's ChatGPT and Codex extensively during the research
process, including for conjecture formation, symbolic experimentation,
proof exploration, consistency checks, and preparation of this manuscript.
The author reviewed the mathematical claims and assumes full responsibility
for the contents.  The theorem uses the exhaustive exact coefficient-scheme
calculation stated in the text and supplied with the paper.

\bibliographystyle{plain}
\bibliography{references}

\end{document}
